% ch8.tex — 练习与补充 IV（中文版）
\chapter{练习与补充 IV}

\section{Sturm--Liouville 问题}

一般的 Sturm--Liouville 本征值问题是
\begin{equation}
-\frac{d}{dx}\!\left(p(x)\frac{dy}{dx}\right) + q(x)\,y = \lambda\,w(x)\,y,
\qquad x\in[a,b],
\label{eq:8.1.1}
\end{equation}
带自伴边界条件
$\alpha_1 y(a)+\alpha_2 y'(a)=0$，
$\beta_1 y(b)+\beta_2 y'(b)=0$。

本征值 $\lambda_n$ 是实数且有序排列
$\lambda_0<\lambda_1<\lambda_2<\cdots$，本征函数正交：
\begin{equation}
\int_a^b y_n(x)\,w(x)\,y_m(x)\,dx = N_n\delta_{nm}.
\label{eq:8.1.2}
\end{equation}

完备性：任何满足边界条件的 $f(x)$ 都可以展开为
$f(x)=\sum_n c_n y_n(x)$，其中
$c_n = \frac{1}{N_n}\int_a^b f(x)w(x)y_n(x)dx$。

\begin{exercise}
对 $[0,\pi]$ 上 $p=1$、$q=0$、$w=1$、带 Dirichlet 边界条件的情形，
求 $\lambda_n$、$y_n(x)$ 并验证正交性。
\end{exercise}

\section{Green 函数}

对于 $\hat L y = s(x)$，解为
\begin{equation}
y(x) = \int_a^b G(x,\xi)\,s(\xi)\,d\xi,
\label{eq:8.2.1}
\end{equation}
其中 Green 函数满足
$(\hat L-\lambda)G(x,\xi) = \delta(x-\xi)$（或 $\lambda=0$ 时 $\hat L G=\delta$）。

按本征函数展开构造：
\begin{equation}
G(x,\xi) = \sum_{n=0}^{\infty}\frac{y_n(x)\,y_n(\xi)}{N_n(\lambda_n-\lambda)}.
\label{eq:8.2.2}
\end{equation}

二阶 ODE 的跳跃条件：
$\left.\frac{\partial G}{\partial x}\right|_{x=\xi^+}
-\left.\frac{\partial G}{\partial x}\right|_{x=\xi^-} = \frac{1}{p(\xi)}$。

\begin{exercise}
对 $[0,1]$ 上 $y''=s(x)$ 且 $y(0)=y(1)=0$ 构造 $G(x,\xi)$，
并验证 $\int_0^1 G(x,\xi)\,d\xi = x(1-x)/2$。
\end{exercise}

\section{Laguerre 多项式、相关算子与 PDE}

Laguerre 导数是
$\hat\Lambda = x\frac{d^2}{dx^2}+(1-x)\frac{d}{dx}$。
它作用于 Laguerre 多项式为
\begin{equation}
\hat\Lambda L_n(x) = -n\,L_n(x).
\label{eq:8.3.1}
\end{equation}

Laguerre 热方程
$\partial_t F = \hat\Lambda F$（$F(x,0)=x^n$）
的解为
\begin{equation}
F(x,t) = L_n\!\left(\frac{x}{1-t}\right)(1-t)^n,
\qquad |t|<1.
\label{eq:8.3.2}
\end{equation}

Umbral 表示：
$L_n(x)=(\hat\theta)_n\,e^{-x\hat\theta}$，其中
$(\hat\theta)_n=\hat\theta(\hat\theta-1)\cdots(\hat\theta-n+1)$
且 $\hat\theta e^{-x\hat\theta}=e^{-x\hat\theta}$。

\begin{exercise}
用 umbral 表示验证 \eqref{eq:8.3.1}，并推导生成函数
$\sum_{n=0}^\infty L_n(x)t^n = e^{-xt/(1-t)}/(1-t)$。
\end{exercise}

\section{Appell 多项式、相关算子与偏微分方程}

Appell 多项式 $A_n(x)$ 满足
$\frac{d}{dx}A_n(x)=n\,A_{n-1}(x)$，并由
\begin{equation}
\sum_{n=0}^{\infty} A_n(x)\,\frac{t^n}{n!}=e^{x t}\,A(t),
\qquad A(0)=1.
\label{eq:8.4.1}
\end{equation}
生成。

满足 $\hat D_A A_n(x)=n\,A_{n-1}(x)$ 的相关求导算子 $\hat D_A$ 通过
$F(x,t)=e^{t\hat D_A}P(x)$ 求解 Cauchy 问题
$\partial_t F=\hat D_A F$，$F(x,0)=P(x)$。

\begin{exercise}
对于 $A_n(x)=x^n$（普通单项式），证明
$\hat D_A=\frac{d}{dx}$，并从 $e^{t\hat D_A}$ 恢复 Taylor 定理。
\end{exercise}

\section{Riemann 函数}

Riemann 方法求解二阶双曲型 PDE 的 Cauchy 问题。对于
\begin{equation}
\frac{\partial^2 u}{\partial x\partial y} + a(x,y)\frac{\partial u}{\partial x}
+ b(x,y)\frac{\partial u}{\partial y} + c(x,y)u = 0,
\label{eq:8.5.1}
\end{equation}
Riemann--Green 函数 $R(x,y;\xi,\eta)$ 满足带特征边界条件的伴随方程。

带非特征曲线上的 Cauchy 数据的解为
\begin{equation}
u(P) = \frac12\left[u(A)R(P;A)+u(B)R(P;B)\right]
+ \frac12\int_{AB}\left(R\frac{\partial u}{\partial n}-u\frac{\partial R}{\partial n}\right)ds.
\label{eq:8.5.2}
\end{equation}

\begin{exercise}
对于一维波动方程 $u_{xy}=0$，证明 $R=1$ 并恢复 d'Alembert 公式。
\end{exercise}

\section{Bessel 特殊函数}

Bessel 函数 $J_n(x)$ 满足
\begin{equation}
x^2J_n''+xJ_n'+(x^2-n^2)J_n=0,
\label{eq:8.6.1}
\end{equation}
生成函数为
\begin{equation}
e^{\frac{x}{2}(t-1/t)} = \sum_{n=-\infty}^{\infty} J_n(x)\,t^n.
\label{eq:8.6.2}
\end{equation}

递推关系：
\begin{align}
J_{n-1}+J_{n+1} &= \frac{2n}{x}J_n,\\
J_{n-1}-J_{n+1} &= 2J_n'.
\label{eq:8.6.3}
\end{align}

有限区间上带权重 $x$ 的正交性：
\begin{equation}
\int_0^a x\,J_n\!\left(\frac{\alpha_{nm}}{a}x\right)
J_n\!\left(\frac{\alpha_{nk}}{a}x\right)dx
= \frac{a^2}{2}J_{n+1}^2(\alpha_{nm})\,\delta_{mk},
\label{eq:8.6.4}
\end{equation}
其中 $\alpha_{nm}$ 是 $J_n$ 的第 $m$ 个零点。

修正 Bessel 函数：$I_n(x)=i^{-n}J_n(ix)$ 满足
$x^2I_n''+xI_n'-(x^2+n^2)I_n=0$。

Hankel 函数：$H_n^{(1)}=J_n+iY_n$，
$H_n^{(2)}=J_n-iY_n$，描述出射/入射柱波。

\begin{exercise}
由生成函数 \eqref{eq:8.6.2} 推导递推关系 \eqref{eq:8.6.3}。
\end{exercise}

\begin{exercise}
利用 $J_0$ 的 Laplace 变换计算 $\int_0^\infty e^{-ax}J_0(bx)\,dx=1/\sqrt{a^2+b^2}$。
\end{exercise}

% 图形
\begin{figure}[h]
\centering
\includegraphics[max width=0.85\linewidth,keepaspectratio]{fig8_1.png}
\caption{Bessel 函数 $J_n(x)$，$n=0,1,2,3$。}
\label{fig:8.1}
\end{figure}

\begin{figure}[h]
\centering
\includegraphics[max width=0.85\linewidth,keepaspectratio]{fig8_3.png}
\caption{波动方程的 Riemann 函数。}
\label{fig:8.3}
\end{figure}
